% ---------------------------------------------------------------------------
% Póster del trabajo final — Visión por Computador (Maestría en IA, Grupo 3)
% Formato: A0 vertical (841 x 1189 mm), en inglés, 2 columnas, UNA sola página.
% Solo paquetes base de LaTeX (sin tikzposter/beamerposter) para que compile
% en cualquier instalación: geometry, multicol, xcolor, graphicx, booktabs.
% Compilar: pdflatex -interaction=nonstopmode poster.tex   (2 pasadas)
% ---------------------------------------------------------------------------
\documentclass{article}
\usepackage[a0paper,portrait,top=2.6cm,bottom=2.2cm,left=2.8cm,right=2.8cm]{geometry}
\usepackage[utf8]{inputenc}
\usepackage[T1]{fontenc}
\usepackage{graphicx}
\usepackage{multicol}
\usepackage{xcolor}
\usepackage{booktabs}
\usepackage{array}
\usepackage{helvet}
\usepackage{url}
\renewcommand{\familydefault}{\sfdefault}
\pagestyle{empty}
\setlength{\parindent}{0pt}
\setlength{\columnsep}{2.8cm}
\setlength{\tabcolsep}{8pt}
\renewcommand{\arraystretch}{1.45}

\definecolor{azul}{HTML}{0F2E52}      % institucional, encabezados
\definecolor{acento}{HTML}{C8102E}    % rojo: cifra principal
\definecolor{verde}{HTML}{1B7F5A}     % verde: hallazgo positivo
\definecolor{gris}{HTML}{EEF1F5}      % fondo de cajas
\definecolor{grisborde}{HTML}{C6CFDA}

% --- tamaños de letra (A0: lectura cómoda a ~1.5 m) ---
\newcommand{\cuerpo}{\fontsize{28}{36}\selectfont}
\newcommand{\chico}{\fontsize{23}{29}\selectfont}
\newcommand{\seccion}[1]{%
  \vspace{11mm}%
  \noindent\colorbox{azul}{\parbox{\dimexpr\columnwidth-2\fboxsep\relax}{%
    \vspace{1.5mm}\color{white}\bfseries\fontsize{37}{44}\selectfont #1\vspace{1.5mm}}}%
  \par\vspace{4mm}}
\newcommand{\caja}[1]{%
  \vspace{2mm}\noindent\fcolorbox{grisborde}{gris}{%
    \parbox{\dimexpr\columnwidth-2\fboxsep-2\fboxrule\relax}{\vspace{3mm}#1\vspace{3mm}}}%
  \par\vspace{2mm}}
\newcommand{\etapa}[2]{%
  \colorbox{azul}{\parbox[c][3.1cm][c]{0.157\columnwidth}{\centering\color{white}%
  \bfseries\fontsize{22}{26}\selectfont #1\\[1pt]{\mdseries\fontsize{16}{19}\selectfont #2}}}}
\newcommand{\flecha}{\raisebox{1.1cm}{\color{azul}\fontsize{30}{30}\selectfont$\rightarrow$}}

\begin{document}
\cuerpo

% =========================== ENCABEZADO ===================================
\noindent
\begin{minipage}[c]{0.115\textwidth}
  
\includegraphics[width=\linewidth]{figuras/logo_uni.png}
\end{minipage}%
\begin{minipage}[c]{0.765\textwidth}
  \centering
  {\bfseries\fontsize{68}{80}\selectfont\color{azul}
  Automatic Extraction of Electoral Results from the\\[4pt]
  Tally Sheets of the 2026 Peruvian General Elections\par}
  \vspace{7mm}
  {\fontsize{29}{36}\selectfont
  Josemanuel Rossy Ca\~nari Palante \quad$\cdot$\quad Kenny Asto Hinostroza
  \quad$\cdot$\quad Melissa Dessire Aylas Barranca \quad$\cdot$\quad Carlos P\'erez P\'erez\par}
  \vspace{4mm}
  {\fontsize{25}{31}\selectfont\itshape
  Master's Program in Artificial Intelligence --- Computer Vision, Group 3 ---
  Universidad Nacional de Ingenier\'ia, Lima, Peru\par}
\end{minipage}%
\begin{minipage}[c]{0.115\textwidth}
  \raggedleft
  
\includegraphics[width=0.78\linewidth]{figuras/qr_repo.png}\\[2pt]
  {\raggedleft\chico\bfseries\color{azul}code + data\par}
\end{minipage}

\vspace{6mm}
{\color{azul}\rule{\textwidth}{5pt}}
\vspace{4mm}

% ==================== FRANJA DE RESULTADOS CLAVE ==========================
\noindent
\fcolorbox{grisborde}{gris}{\parbox{\dimexpr\textwidth-2\fboxsep-2\fboxrule\relax}{%
\vspace{4mm}
\begin{minipage}[t]{0.315\linewidth}\centering
  {\bfseries\fontsize{80}{88}\selectfont\color{acento}59.6\%}\\[4pt]
  {\fontsize{25}{31}\selectfont field accuracy on handwritten actas\\
  95\% CI [56.6, 62.5] $\cdot$ CER 0.401 $\cdot$ 3{,}268 fields}
\end{minipage}\hfill
\begin{minipage}[t]{0.315\linewidth}\centering
  {\bfseries\fontsize{80}{88}\selectfont\color{verde}95.1\%}\\[4pt]
  {\fontsize{25}{31}\selectfont on born-digital actas, parsed from the PDF text
  layer \textbf{with no vision at all}}
\end{minipage}\hfill
\begin{minipage}[t]{0.315\linewidth}\centering
  {\bfseries\fontsize{80}{88}\selectfont\color{azul}0}\\[4pt]
  {\fontsize{25}{31}\selectfont manually annotated fields: ground truth is the
  official result published by the same portal}
\end{minipage}
\vspace{4mm}}}

\vspace{7mm}

% =========================== COLUMNAS ======================================
\begin{multicols}{2}

\seccion{1. Problem}
Peru's National Office of Electoral Processes (ONPE) publishes every tally
sheet (\emph{acta de escrutinio}) digitally, but the information is
\textbf{visual}: a scanned form where polling-station members wrote the vote
counts by hand. Manual extraction does not scale to the ${\sim}88{,}000$
stations of a general election, which limits automatic auditing and
independent verification. \textbf{Our question:} how far can a classical
computer vision pipeline go on \emph{real, nationally sampled} actas, and
what does that imply for a production system?

\seccion{2. Dataset and ground truth}
\textbf{Source.} Official ONPE portal, 2026 General Elections, first round
(presidential); actas downloaded through the portal's own API.

\vspace{3mm}
\textbf{Sampling.} Simple random sample without replacement over the
${\sim}88{,}064$ stations with a \textbf{published seed (2026)}: 100 actas,
100\% hit rate, 23 department-level codes, size-proportional (Lima 36\%).
This removes the bias of the pilot, which used sequential codes from a single
district.

\vspace{3mm}
\textbf{Annotation-free ground truth.} The same portal publishes the
officially digitized counts, so every acta yields ${\sim}43$ labeled
field--value pairs \textbf{at zero annotation cost} and the evaluation is
fully reproducible.

\caja{\textbf{Population finding with architectural consequences.} The
population is \emph{mixed}: \textbf{76 handwritten} scans and \textbf{24
born-digital (STAE)} documents with an extractable text layer (prevalence
24\%, 95\% Wilson CI [16.7, 33.2]). A production system must
\textbf{classify the document type first and route accordingly}.}

\seccion{3. Method}
\vspace{-1mm}
\begin{center}
\etapa{1. Acquire}{API + PDF$\rightarrow$PNG 300 dpi}\,\flecha\,
\etapa{2. Preprocess}{deskew, CLAHE, denoise}\,\flecha\,
\etapa{3. Detect}{registered template}\,\flecha\,
\etapa{4. Recognize}{digit-restricted OCR}\,\flecha\,
\etapa{5. Structure}{JSON per acta}
\end{center}
\vspace{2mm}

\textbf{Detection by fiducial registration.} The acta is a fixed-layout form,
so 46 regions (38 party rows, 4 summary rows, 4 context fields) are defined
as page fractions. Real scans break fixed fractions (black bands, offsets),
so the 15 printed registration squares are detected, matched to a reference
acta by \textbf{RANSAC on translation}, and corrected with a partial affine
warp before cropping.

\vspace{3mm}
\textbf{Recognition.} EasyOCR restricted to digits over cleaned cells:
interior margin against printed borders, ink isolation by dark threshold and
Otsu, erasure of the dotted separators by template position, speckle
filtering.

\vspace{3mm}
\textbf{Domain rule.} Polling-station members leave a cell \emph{blank}
instead of writing 0, so an empty cleaned cell is read as \textbf{0 votes} ---
a semantic decision, not a vision one.

\vspace{4mm}
\begin{center}
\includegraphics[height=28cm]{figuras/plantilla_overlay.png}
\hspace{1.2cm}
\includegraphics[height=28cm]{figuras/limpieza_celdas.png}\\[2mm]
{\chico Left: template regions after fiducial registration. Right: cell
cleaning (raw crop $\rightarrow$ ink isolation $\rightarrow$ separators erased).}
\end{center}

\columnbreak

\seccion{4. Experimental setup}
\textbf{Metrics.} \emph{Field accuracy} (exact match against the official
integer) and \emph{CER} over digit strings.

\vspace{3mm}
\textbf{Statistics.} Fields inside one acta share writer and scan quality, so
\textbf{the acta is the cluster}: every interval comes from a
\textbf{clustered bootstrap} (10,000 resamples of whole actas). Per-field
binomial intervals would be optimistic --- the measured design effect is
\textbf{3.07} on handwritten actas (76 actas $\approx$ 1,066 independent
fields). Ablations are re-evaluated on the saved outputs, with no OCR re-run.

\seccion{5. Results}
\textbf{National random sample --- 76 handwritten actas, 3,268 fields}
\vspace{2mm}

{\chico
\begin{tabular}{@{}p{0.62\columnwidth}rr@{}}
\toprule
\textbf{Configuration} & \textbf{Acc.\ (\%)} & \textbf{CER} \\
\midrule
Fixed template, raw OCR            & 26.81 & 0.724 \\
Fixed template + zero rule         & 50.86 & 0.482 \\
Registered (fiducials), raw OCR    & 38.00 & 0.617 \\
\textbf{Registered + zero rule}    & \textbf{59.58} & \textbf{0.401} \\
\bottomrule
\end{tabular}}

\vspace{4mm}
Fiducial registration is worth \textbf{+8.7 points} and lifts the worst acta
from 14\% to 30\%; the domain rule is worth \textbf{+21.6 points}. Per-acta
accuracy averages 59.6\% (SD 13.0, range 30--84\%): the residual variance is
handwriting and scan quality, not misalignment.

\vspace{5mm}
\textbf{Pilot ablation --- 10 actas, 430 fields}
\vspace{2mm}

{\chico
\begin{tabular}{@{}p{0.62\columnwidth}rr@{}}
\toprule
\textbf{Configuration} & \textbf{Acc.\ (\%)} & \textbf{CER} \\
\midrule
v1: raw digit OCR                        & 43.95 & 0.524 \\
v2: + ink isolation, borders, sub-boxes  & 52.09 & 0.419 \\
v3: + separator erasure, Otsu, speckle   & 52.79 & 0.476 \\
v4: + empty-cell-is-zero rule            & 58.37 & 0.420 \\
\bottomrule
\end{tabular}}

\vspace{5mm}
\textbf{Does a multimodal LLM replace the pipeline? No.} Same 15 handwritten
actas (645 fields), same protocol; the LLM gets the full page plus two
image--JSON exemplars, with no template and no OCR.
\vspace{2mm}

{\chico
\begin{tabular}{@{}p{0.62\columnwidth}rr@{}}
\toprule
\textbf{Method} & \textbf{Acc.\ (\%)} & \textbf{CER} \\
\midrule
\textbf{Pipeline (registered template + OCR)} & \textbf{60.93} & \textbf{0.395} \\
LLM few-shot, end-to-end (2 exemplars)        & 48.37 & 0.591 \\
\bottomrule
\end{tabular}}

\vspace{4mm}
The LLM reads the layout well but is not trained for handwritten Peruvian
digits and collapses on poor scans (20.9\% and 23.3\% on the two worst actas,
where the pipeline still reads 34.9\% and 53.5\%). Its real advantage is
\textbf{zero detection engineering}: a bootstrap pre-labeler for a new form,
not a replacement recognizer.

\caja{\textbf{Is 24 actas enough for the 95.1\%?} Yes --- text parsing is
deterministic, so between-acta variance falls below binomial (ICC
$\approx 0$, design effect 1.0) and the interval is [94.5, 95.7], nine times
tighter than the handwritten one. The 35.5-point gap has a 95\% CI of
[32.6, 38.5].}

\seccion{6. Conclusions and future work}
\textbf{1.} Register the template on the form's fiducial marks \emph{before}
any fixed-coordinate cropping.\\[2mm]
\textbf{2.} Separate domain semantics (blank $=$ 0) from vision performance;
conflating them hides where the system really fails.\\[2mm]
\textbf{3.} Classify the document type first: routing born-digital actas to
text parsing is worth \textbf{35.5 points}.\\[2mm]
\textbf{4.} Generality does not yet beat specialization: a general multimodal
LLM is cheaper to build but 12.6 points less accurate.

\caja{\textbf{Limitations.} 59.6\% is a \emph{ceiling for generic OCR}, not a
deployable accuracy. Errors now concentrate in \textbf{recognition}, not
localization --- exactly the component fine-tuning replaces. Next: fine-tune
TrOCR on the free field--value pairs (${\sim}43{,}000$ from 1,000 actas),
train YOLOv11 / RT-DETR on labels auto-generated by the registered template,
and use the LLM as pre-labeler in the human-correction loop.}

\end{multicols}

\vspace{3mm}
{\color{azul}\rule{\textwidth}{4pt}}
\vspace{2mm}

\begin{center}
{\chico Code, per-field evaluation artifacts and the sampling seed:
\textbf{\url{https://github.com/carlosperez100/onpe_actas}} \quad$\cdot$\quad
Related work: RT-DETR (CVPR 2024) $\cdot$ DTrOCR (WACV 2024) $\cdot$
OmniParser (CVPR 2024)}
\end{center}

\end{document}
